% Generated draft: requires mathematical and source review.
\documentclass[11pt]{article}
\usepackage[margin=25mm]{geometry}
\usepackage{fontspec}
\setmainfont{Latin Modern Roman}
\usepackage{amsmath,amssymb,mathtools,unicode-math}
\setmathfont{Latin Modern Math}
\usepackage{enumitem,needspace,xurl,hyperref}
\setlength{\parindent}{0pt}
\setlength{\parskip}{5pt}
\setlength{\emergencystretch}{3em}
\newcommand{\N}{\mathbb{N}}
\newcommand{\Z}{\mathbb{Z}}
\newcommand{\Q}{\mathbb{Q}}
\newcommand{\R}{\mathbb{R}}
\newcommand{\C}{\mathbb{C}}
\newcommand{\F}{\mathbb{F}}
\newcommand{\A}{\mathbb{A}}
\newcommand{\PP}{\mathbb{P}}
\newcommand{\E}{\mathbb{E}}
\newcommand{\eps}{\varepsilon}
\newcommand{\dd}{\,\mathrm{d}}
\begin{document}
\section*{UnsolvedMath ID: 20003315}
\textbf{Source problem code:} catalogue entry
\subsection{Definitions and mathematical statement}
Let \(X_3=\operatorname{SL}_3(\R)/\operatorname{SL}_3(\Z)\), and let \(A\) be the full diagonal subgroup. Its positive identity component is
\[
 A^+=\{\operatorname{diag}(e^{u},e^{v},e^{-u-v}):u,v\in\R\}.
\]
It has finite index in \(A\), so boundedness of the two full-group orbits is equivalent. Bounded means relatively compact in \(X_3\).

A nontrivial diagonal ray is
\[
 A_r^+=\{\operatorname{diag}(e^{tr_1},e^{tr_2},e^{tr_3}):t\geq0\},
 \quad r_1+r_2+r_3=0,\quad r\ne0.
\]
The recovered source conjecture asks: for every prescribed pair of such rays, is there \(x\in X_3\) with both ray orbits bounded but \(Ax\) unbounded? The extracted note about rays in the cone \(u,v\geq0\) ends mid-sentence; it is not completed or made an extra restriction on the main question.

The aid reports the stronger statement that for every finite or countable prescribed family \((A_{r_j}^+)\), the set
\[
 E=\{x\in X_3:A_{r_j}^+x\text{ is bounded for every }j,\ Ax\text{ is unbounded}\}
\]
is thick: for every nonempty open \(U\subset X_3\),
\[
 \dim_{\mathrm H}(E\cap U)=8=\dim X_3.
\]
Hausdorff dimension is taken in a smooth Riemannian metric, with the usual local invariance under smooth coordinate changes. The compact bound for a ray may depend on \(j\) and \(x\); no uniform compact set for the whole countable family is asserted. The result is not quantified over every uncountable family of directions.

The supplied proof route is to let \(F_j\) be the two-sided one-parameter diagonal subgroup containing the \(j\)-th ray. The aid invokes An--Guan--Kleinbock's 2015 full-dimension theorem to say that
\[
 \bigcap_j\{x:F_jx\text{ is bounded}\}
\]
is thick. This is a stronger sufficient boundedness condition than ray boundedness, not an asserted equivalence. It then invokes Einsiedler--Katok--Lindenstrauss, Theorem 10.2, for the Hausdorff-dimension-two bound on
\[
 B_A=\{x:A^+x\text{ is bounded}\}.
\]
Removing that lower-dimensional set leaves a thick subset of \(E\). The argument does not classify all bounded full diagonal orbits or claim that they are all compact closed orbits. The exact hypotheses and full references of both literature inputs require review.

Taking two arbitrary ray directions gives the aid's affirmative answer to the recovered Conjecture 33, including directions not confined to the displayed positive cone. This is a supplied solved-in-literature assessment of that recovered statement; the unfinished source note remains unresolved as text.
\subsection{Short English statement}
The supplied literature deduction gives a thick set of lattices bounded along any prescribed countable collection of nontrivial diagonal rays, yet unbounded under the whole diagonal group. Each ray is bounded separately. The conclusion uses two cited dimension theorems and does not classify the bounded full-group orbits.
\subsection{Sources}
\begin{itemize}
\item Local UnsolvedMath registry, permanent ID \(20003315\), exact source code \texttt{catalogue entry}; canonical \texttt{aim-other-notes.json}, \texttt{notes[121]}, AIM domain Other, registry category Miscellaneous, source tag \texttt{conjecture}.
\item AIM workshop ``Emerging applications of measure rigidity,'' Conjecture 33: \url{https://aimath.org/WWN/measrigid/measrigid.pdf}. Statement status \texttt{reconstructed\_unverified}; the stored note ends after the cone \(u,v\geq0\). The supplied verification refers to official-PDF wording, but no PDF was fetched here.
\item An--Guan--Kleinbock, 2015 full-dimension theorem, and Einsiedler--Katok--Lindenstrauss, Theorem 10.2, as invoked by the aid. No paper titles, identifiers, or full authoritative citations are supplied. The original literature field is empty.
\item Research attempt 1: result \texttt{full\_solution}, problem status at run \texttt{solved}, rigor label \texttt{proved}, novelty confidence \texttt{low}; classification \texttt{SOLVED-IN-LITERATURE}, registry status \texttt{solved}. Created and updated 14 August 2026. These are supplied assessments for the recovered statement, not independent verification.
\item Attribution: Ulam AI / UnsolvedMath Contributors, CC BY 4.0. The machine-generated research aid is not peer reviewed and requires independent verification before citation.
\end{itemize}
\end{document}
